\documentclass[11pt]{amsart}

\usepackage[margin=2.5cm]{geometry}
\usepackage{amsmath,amssymb,amsthm,mathtools}
\usepackage{hyperref}

\newtheorem{theorem}{Theorem}[section]
\newtheorem{proposition}[theorem]{Proposition}
\newtheorem{lemma}[theorem]{Lemma}
\theoremstyle{remark}
\newtheorem{remark}[theorem]{Remark}

\DeclareMathOperator{\Var}{Var}
\DeclareMathOperator{\dist}{dist}
\DeclareMathOperator{\spanop}{span}
\DeclareMathOperator{\diag}{diag}

\newcommand{\R}{\mathbb R}
\newcommand{\cO}{\mathcal O}
\newcommand{\BL}{\mathrm{BL}}

\title[A counterexample to uniform $L^2$ stability]{A counterexample to uniform $L^2$ stability\\ in the Brascamp--Lieb inequality}
\author{}
\date{}

\begin{document}

\begin{abstract}
We write down a short paper whose mathematical content matches the Lean
development in the directory \texttt{Lean/}. The argument produces an explicit
one-dimensional family $(\phi_S,f_S)$ for which
\[
  \frac{\dist_{L^2(\rho_{\phi_S})}(f_S,\cO_{\BL,\phi_S})^2}
       {\delta_{\BL,\phi_S}(f_S)}
  = S\bigl(1+O(S^{-1})\bigr),
  \qquad S\to\infty,
\]
where $\delta_{\BL,\phi}$ denotes the Brascamp--Lieb deficit and
$\cO_{\BL,\phi}=\spanop\{1,\phi'\}$ is the optimizer space in dimension one.
In particular no finite $L^2$-stability constant can hold uniformly along this
family. We then lift the example to every dimension $d\ge 1$ by the product
potential
\[
  \Phi_{S,d}(x,y)=\phi_S(x)+\frac{|y|^2}{2},
\]
and obtain the exact higher-dimensional statement encoded in the Lean theorem:
for each fixed $d\ge 1$ there is no finite constant $C$ such that
\[
  \dist_{L^2(\rho_{\Phi_{S,d}})}(F_{S,d},\cO_{\BL,\Phi_{S,d}})^2
  \le C^2\,\delta_{\BL,\Phi_{S,d}}(F_{S,d})
\]
for every parameter $S$ for which the deficit is positive.
\end{abstract}

\maketitle

\section{Statement}

Let $\phi:\R^d\to\R$ be smooth and strictly convex, with
$0<\int_{\R^d}e^{-\phi(x)}\,dx<\infty$. We write
\[
  Z_\phi:=\int_{\R^d}e^{-\phi(x)}\,dx,
  \qquad
  d\rho_\phi:=Z_\phi^{-1}e^{-\phi(x)}\,dx.
\]
For a locally Lipschitz $f\in L^2(\rho_\phi)$, the Brascamp--Lieb energy and
deficit are
\[
  \mathcal E_\phi(f)
  :=\int_{\R^d}\langle (D^2\phi)^{-1}\nabla f,\nabla f\rangle\,d\rho_\phi,
  \qquad
  \delta_{\BL,\phi}(f):=\mathcal E_\phi(f)-\Var_{\rho_\phi}(f).
\]
The optimizer space is
\[
  \cO_{\BL,\phi}
  :=
  \{a\cdot\nabla\phi+b:\ a\in\R^d,\ b\in\R\}.
\]

The concrete construction below yields a one-dimensional family
$(\phi_S,f_S)$ with $\delta_{\BL,\phi_S}(f_S)>0$ for all sufficiently large
$S$ and
\[
  \frac{\dist_{L^2(\rho_{\phi_S})}(f_S,\cO_{\BL,\phi_S})^2}
       {\delta_{\BL,\phi_S}(f_S)}
  =S\bigl(1+O(S^{-1})\bigr).
\]
This already implies that any stability statement of the form
\[
  \dist_{L^2(\rho_\phi)}(f,\cO_{\BL,\phi})
  \le C\,\delta_{\BL,\phi}(f)^{1/2}
\]
must fail on every class of admissible one-dimensional data that contains the
family $(\phi_S,f_S)$.

The Lean build also proves the following exact fixed-dimension product-family
statement.

\begin{theorem}\label{thm:main}
For every integer $d\ge 1$ there exists a smooth strictly convex family
$\Phi_{S,d}:\R^d\to\R$ and a locally Lipschitz family
$F_{S,d}\in L^2(\rho_{\Phi_{S,d}})$ such that:
\begin{enumerate}
\item for all sufficiently large $S$ one has
  $\delta_{\BL,\Phi_{S,d}}(F_{S,d})>0$;
\item for every finite $C\in\R$ there exists $S$ with
  $\delta_{\BL,\Phi_{S,d}}(F_{S,d})>0$ and
  \[
    \dist_{L^2(\rho_{\Phi_{S,d}})}(F_{S,d},\cO_{\BL,\Phi_{S,d}})^2
    > C^2\,\delta_{\BL,\Phi_{S,d}}(F_{S,d}).
  \]
\end{enumerate}
Equivalently, for fixed $d\ge 1$ there is no finite constant $C$ such that
\[
  \dist_{L^2(\rho_{\Phi_{S,d}})}(F_{S,d},\cO_{\BL,\Phi_{S,d}})^2
  \le C^2\,\delta_{\BL,\Phi_{S,d}}(F_{S,d})
\]
for every parameter $S$ with $\delta_{\BL,\Phi_{S,d}}(F_{S,d})>0$.
\end{theorem}

The proof is one-dimensional at heart. The higher-dimensional statement is
obtained by tensoring with a standard Gaussian in the transverse variables.

\section{The one-dimensional potential}

Fix a non-negative even $C^\infty$ bump
$\kappa:\R\to[0,\infty)$ supported in $[-1,1]$ and normalized by
\[
  \int_{\R}\kappa(u)\,du=1.
\]
For parameters $S\ge 1$, $\varepsilon\in(0,\tfrac14)$ and $\eta>0$, define
\begin{equation}\label{eq:phi2}
  \phi''_{S,\varepsilon,\eta}(x)
  :=
  \eta
  +\frac{S}{\varepsilon}\kappa\!\left(\frac{x-1}{\varepsilon}\right)
  +\frac{S}{\varepsilon}\kappa\!\left(\frac{x+1}{\varepsilon}\right),
\end{equation}
and impose the normalization
\[
  \phi_{S,\varepsilon,\eta}(0)=0,
  \qquad
  \phi'_{S,\varepsilon,\eta}(0)=0.
\]
The right-hand side of \eqref{eq:phi2} is smooth, even, and bounded below by
$\eta$, so $\phi_{S,\varepsilon,\eta}$ is smooth, even, and strictly convex.

We make the specific choice
\begin{equation}\label{eq:choice}
  \varepsilon=S^{-3},
  \qquad
  \eta=S^{-4},
\end{equation}
and abbreviate
\[
  \phi_S:=\phi_{S,S^{-3},S^{-4}},
  \qquad
  \rho_S:=\rho_{\phi_S},
  \qquad
  Z_S:=\int_{\R}e^{-\phi_S(x)}\,dx.
\]
All asymptotic notation below refers to the limit $S\to\infty$.

We decompose $\R$ into the central core, two transition layers, and the
exterior:
\[
  C_S:=[-1+\varepsilon,1-\varepsilon],
  \qquad
  I_S^+:=(1-\varepsilon,1+\varepsilon),
  \qquad
  I_S^-:=(-1-\varepsilon,-1+\varepsilon),
\]
\[
  T_S:=I_S^+\cup I_S^-,
  \qquad
  E_S:=\R\setminus(C_S\cup T_S).
\]

The next lemma summarizes the explicit shape of $\phi_S$ needed in the sequel.

\begin{lemma}\label{lem:profile}
For the choice \eqref{eq:choice}, the following hold for all sufficiently large
$S$:
\begin{enumerate}
\item on $C_S$ one has
  \[
    \phi''_S\equiv\eta,
    \qquad
    \phi'_S(x)=\eta x,
    \qquad
    \phi_S(x)=\frac{\eta x^2}{2}\in\left[0,\frac{\eta}{2}\right];
  \]
\item on $T_S$ one has $\phi_S(x)=O(S^{-2})$ uniformly, hence
  $e^{\phi_S(x)}=1+O(S^{-2})$ uniformly on $T_S$;
\item for $x>1+\varepsilon$,
  \[
    \phi'_S(x)=S+\eta x,
  \]
  and therefore
  \[
    \phi_S(x)
    =\phi_S(1+\varepsilon)
     +S(x-1-\varepsilon)
     +\frac{\eta}{2}\bigl(x^2-(1+\varepsilon)^2\bigr);
  \]
\item by evenness, the corresponding formulas on the left tail follow from the
  right-tail ones by the substitution $x\mapsto -x$.
\end{enumerate}
\end{lemma}

\begin{proof}
On the core $C_S$ both translated bumps in \eqref{eq:phi2} vanish, so
$\phi''_S=\eta$ there. Integrating and using $\phi'_S(0)=\phi_S(0)=0$ gives
the displayed core formulas.

On $I_S^+$ we have
\[
  \phi_S(x)
  =\phi_S(1-\varepsilon)+\int_{1-\varepsilon}^x\phi'_S(t)\,dt.
\]
Since $\phi_S(1-\varepsilon)\le \eta/2$ and
$\sup_{I_S^+}\phi'_S\le \eta+S$, we obtain
\[
  \phi_S(x)\le \frac{\eta}{2}+2\varepsilon(\eta+S)=O(S^{-2}),
\]
using \eqref{eq:choice}. The same estimate holds on $I_S^-$ by evenness. The
bound on $e^{\phi_S}-1$ is immediate.

On $(1+\varepsilon,\infty)$ both bumps have again vanished, so
$\phi''_S=\eta$. Since the total mass of the right bump is $S$,
$\phi'_S(1+\varepsilon)=S+\eta(1+\varepsilon)$, and the right-tail formulas
follow by integration. The left tail is obtained by symmetry.
\end{proof}

\section{Normalization and the tail mass}

The next lemma is the only normalization input needed later on.

\begin{lemma}\label{lem:normalization}
With the choice \eqref{eq:choice}, as $S\to\infty$:
\begin{enumerate}
\item
  \[
    \int_{1+\varepsilon}^{\infty}e^{-\phi_S(x)}\,dx
    =\frac{1}{S}+O(S^{-3});
  \]
\item
  \[
    Z_S=2+\frac{2}{S}+O(S^{-3});
  \]
\item if
  \[
    q_S:=\rho_S(E_S)=\rho_S(\{|x|>1+\varepsilon\}),
    \qquad
    t_S:=\rho_S(T_S),
  \]
  then
  \[
    q_S=\frac{1}{S}-\frac{1}{S^2}+O(S^{-3}),
    \qquad
    t_S=O(S^{-3}).
  \]
\end{enumerate}
\end{lemma}

\begin{proof}
By Lemma~\ref{lem:profile}, for $x>1+\varepsilon$ we have
\[
  \phi_S(x)
  =\phi_S(1+\varepsilon)+\widetilde S\,(x-1-\varepsilon)
   +\frac{\eta}{2}(x-1-\varepsilon)^2,
\]
where $\widetilde S=S+\eta(1+\varepsilon)=S+O(S^{-4})$. Thus
\[
  \int_{1+\varepsilon}^{\infty}e^{-\phi_S(x)}\,dx
  =
  e^{-\phi_S(1+\varepsilon)}
  \int_0^\infty e^{-\widetilde S u-\eta u^2/2}\,du.
\]
Since $\phi_S(1+\varepsilon)=O(S^{-2})$, the prefactor equals
$1+O(S^{-2})$. Also
\[
  0\le \frac{1}{\widetilde S}
  -\int_0^\infty e^{-\widetilde S u-\eta u^2/2}\,du
  \le
  \frac{\eta}{2}\int_0^\infty u^2e^{-\widetilde S u}\,du
  =
  \frac{\eta}{\widetilde S^3}
  =O(S^{-7}),
\]
so the tail integral is $1/S+O(S^{-6})$. This proves (1).

For (2), split $\int_0^\infty e^{-\phi_S}\,dx$ into the core, the right layer,
and the right tail. On the core, $\phi_S(x)=\eta x^2/2$, hence
\[
  \int_0^{1-\varepsilon}e^{-\phi_S(x)}\,dx
  =1-\varepsilon+O(\eta).
\]
On the layer $I_S^+$, Lemma~\ref{lem:profile} gives
$e^{-\phi_S(x)}=1+O(S^{-2})$ uniformly, so
\[
  \int_{1-\varepsilon}^{1+\varepsilon}e^{-\phi_S(x)}\,dx
  =2\varepsilon+O(S^{-5}).
\]
On the tail we use (1). Therefore
\[
  \int_0^\infty e^{-\phi_S(x)}\,dx
  =1+\frac{1}{S}+O(S^{-3}),
\]
and by evenness
\[
  Z_S=2+\frac{2}{S}+O(S^{-3}).
\]

Finally,
\[
  t_S
  \le \frac{|T_S|}{Z_S}\sup_{x\in T_S}e^{-\phi_S(x)}
  =O(\varepsilon)
  =O(S^{-3}).
\]
Moreover,
\[
  q_S
  =\frac{2}{Z_S}\int_{1+\varepsilon}^\infty e^{-\phi_S(x)}\,dx
  =
  \frac{2/S+O(S^{-3})}{2+2/S+O(S^{-3})}
  =
  \frac{1}{S}-\frac{1}{S^2}+O(S^{-3}),
\]
which proves (3).
\end{proof}

\section{The transition test function}

Define
\begin{equation}\label{eq:g}
  g_S(x):=
  \begin{cases}
    0,& x\in C_S,\\[4pt]
    \displaystyle
    \frac{\int_{1-\varepsilon}^{|x|}\phi''_S(t)e^{\phi_S(t)}\,dt}
         {\int_{I_S^+}\phi''_S(t)e^{\phi_S(t)}\,dt},
      & |x|\in(1-\varepsilon,1+\varepsilon),\\[12pt]
    1,& x\in E_S.
  \end{cases}
\end{equation}
This function is even, continuous, bounded between $0$ and $1$, and globally
Lipschitz.

Set
\[
  A_S:=\int_{I_S^+}\phi''_S(t)e^{\phi_S(t)}\,dt.
\]
The energy of $g_S$ is explicit.

\begin{lemma}\label{lem:energy}
As $S\to\infty$,
\[
  A_S=S\bigl(1+O(S^{-2})\bigr),
\]
and
\[
  \mathcal E_{\phi_S}(g_S)
  =
  \frac{2}{Z_SA_S}
  =
  \frac{1}{S}-\frac{1}{S^2}+O(S^{-3}).
\]
\end{lemma}

\begin{proof}
By Lemma~\ref{lem:profile}, $e^{\phi_S}=1+O(S^{-2})$ uniformly on $I_S^+$,
while
\[
  \int_{I_S^+}\phi''_S(t)\,dt=S+2\eta\varepsilon=S+O(S^{-7}).
\]
Hence
\[
  A_S
  =
  \int_{I_S^+}\phi''_S(t)e^{\phi_S(t)}\,dt
  =
  S\bigl(1+O(S^{-2})\bigr).
\]

To compute the energy, note that $g_S$ is constant on $C_S\cup E_S$, so only
the two layers contribute. On $I_S^+$ one has
\[
  g'_S(x)=\frac{\phi''_S(x)e^{\phi_S(x)}}{A_S},
\]
and the same formula holds on $I_S^-$ up to the sign change forced by
evenness. Therefore
\[
  \mathcal E_{\phi_S}(g_S)
  =
  \frac{2}{Z_S}\int_{I_S^+}\frac{(g'_S(x))^2}{\phi''_S(x)}e^{-\phi_S(x)}\,dx
  =
  \frac{2}{Z_SA_S}.
\]
Combining this with Lemma~\ref{lem:normalization} gives
\[
  \frac{2}{Z_SA_S}
  =
  \frac{2}{(2+2/S+O(S^{-3}))\,S(1+O(S^{-2}))}
  =
  \frac{1}{S}-\frac{1}{S^2}+O(S^{-3}).
\]
\end{proof}

\begin{lemma}\label{lem:variance}
The variance of $g_S$ satisfies
\[
  \Var_{\rho_S}(g_S)
  =
  q_S(1-q_S)+O(S^{-3})
  =
  \frac{1}{S}-\frac{2}{S^2}+O(S^{-3}).
\]
\end{lemma}

\begin{proof}
Since $g_S=0$ on $C_S$, $g_S=1$ on $E_S$, and $0\le g_S\le 1$ on $T_S$, we
have
\[
  \int g_S\,d\rho_S=q_S+O(t_S),
  \qquad
  \int g_S^2\,d\rho_S=q_S+O(t_S).
\]
By Lemma~\ref{lem:normalization}, $t_S=O(S^{-3})$, so
\[
  \Var_{\rho_S}(g_S)
  =
  \int g_S^2\,d\rho_S-\left(\int g_S\,d\rho_S\right)^2
  =
  q_S-q_S^2+O(S^{-3}).
\]
Substituting the expansion for $q_S$ from
Lemma~\ref{lem:normalization} gives the result.
\end{proof}

Now define
\[
  c_S:=\int_{\R}g_S\,d\rho_S,
  \qquad
  f_S:=g_S-c_S.
\]
Then $f_S$ is even and has $\rho_S$-mean zero. Since the deficit and the
variance are invariant under addition of constants,
\begin{equation}\label{eq:constant-shift}
  \mathcal E_{\phi_S}(f_S)=\mathcal E_{\phi_S}(g_S),
  \qquad
  \Var_{\rho_S}(f_S)=\Var_{\rho_S}(g_S),
  \qquad
  \delta_{\BL,\phi_S}(f_S)=\delta_{\BL,\phi_S}(g_S).
\end{equation}

\begin{lemma}\label{lem:orth}
The function $f_S$ is orthogonal in $L^2(\rho_S)$ to
\[
  \cO_{\BL,\phi_S}=\spanop\{1,\phi'_S\}.
\]
Consequently
\[
  \dist_{L^2(\rho_S)}(f_S,\cO_{\BL,\phi_S})^2
  =
  \|f_S\|_{L^2(\rho_S)}^2
  =
  \Var_{\rho_S}(g_S).
\]
\end{lemma}

\begin{proof}
By construction $\int f_S\,d\rho_S=0$, so $f_S\perp 1$. Also $f_S$ is even,
$\phi'_S$ is odd, and $e^{-\phi_S}$ is even, hence
$f_S(x)\phi'_S(x)e^{-\phi_S(x)}$ is odd and integrable. Therefore
\[
  \int_{\R}f_S(x)\phi'_S(x)\,d\rho_S(x)=0.
\]
Thus $f_S\perp \spanop\{1,\phi'_S\}$. The distance identity follows because an
orthogonal projection onto $\cO_{\BL,\phi_S}$ vanishes on $f_S$, while
\[
  \|f_S\|_{L^2(\rho_S)}^2=\Var_{\rho_S}(g_S)
\]
since $f_S$ is the centered version of $g_S$.
\end{proof}

\begin{proposition}\label{prop:asymptotics}
As $S\to\infty$,
\[
  \mathcal E_{\phi_S}(f_S)=\frac{1}{S}-\frac{1}{S^2}+O(S^{-3}),
  \qquad
  \Var_{\rho_S}(f_S)=\frac{1}{S}-\frac{2}{S^2}+O(S^{-3}),
  \qquad
  \delta_{\BL,\phi_S}(f_S)=\frac{1}{S^2}+O(S^{-3}).
\]
In particular, $\delta_{\BL,\phi_S}(f_S)>0$ for all sufficiently large $S$.
\end{proposition}

\begin{proof}
The first identity follows from Lemma~\ref{lem:energy} and
\eqref{eq:constant-shift}. The second follows from Lemma~\ref{lem:variance}
and \eqref{eq:constant-shift}. Subtracting gives
\[
  \delta_{\BL,\phi_S}(f_S)
  =
  \left(\frac{1}{S}-\frac{1}{S^2}\right)
  -
  \left(\frac{1}{S}-\frac{2}{S^2}\right)
  +O(S^{-3})
  =
  \frac{1}{S^2}+O(S^{-3}),
\]
and the eventual positivity is immediate.
\end{proof}

\section{The one-dimensional contradiction}

\begin{theorem}\label{thm:1d}
The family $(\phi_S,f_S)$ satisfies
\[
  \frac{\dist_{L^2(\rho_S)}(f_S,\cO_{\BL,\phi_S})^2}
       {\delta_{\BL,\phi_S}(f_S)}
  =
  S\bigl(1+O(S^{-1})\bigr),
  \qquad S\to\infty.
\]
In particular, for every finite $C$ there exists $S$ such that
$\delta_{\BL,\phi_S}(f_S)>0$ and
\[
  \dist_{L^2(\rho_S)}(f_S,\cO_{\BL,\phi_S})^2
  >
  C^2\,\delta_{\BL,\phi_S}(f_S).
\]
\end{theorem}

\begin{proof}
By Lemma~\ref{lem:orth} and Proposition~\ref{prop:asymptotics},
\[
  \dist_{L^2(\rho_S)}(f_S,\cO_{\BL,\phi_S})^2
  =
  \frac{1}{S}-\frac{2}{S^2}+O(S^{-3}),
\]
while
\[
  \delta_{\BL,\phi_S}(f_S)=\frac{1}{S^2}+O(S^{-3}).
\]
Therefore
\[
  \frac{\dist^2}{\delta_{\BL}}
  =
  \frac{(1/S)\bigl(1-2/S+O(S^{-2})\bigr)}
       {(1/S^2)\bigl(1+O(S^{-1})\bigr)}
  =
  S\bigl(1+O(S^{-1})\bigr).
\]
The second assertion follows because the ratio tends to $+\infty$ and the
deficit is eventually positive.
\end{proof}

\begin{remark}\label{rem:lean-1d}
The Lean theorem in dimension one is phrased in a slightly more abstract way:
it records only the numerical invariants ``distance squared'' and ``deficit''
of an admissible datum. Theorem~\ref{thm:1d} is exactly the concrete
counterexample needed for that abstract formulation: any class of admissible
one-dimensional Brascamp--Lieb data that contains the family $(\phi_S,f_S)$
cannot satisfy a uniform square-root stability estimate.
\end{remark}

\section{The higher-dimensional lift}

Fix $d\ge 2$. Write points of $\R^d$ as $(x,y)\in \R\times\R^{d-1}$ and define
\[
  \Phi_{S,d}(x,y):=\phi_S(x)+\frac{|y|^2}{2},
  \qquad
  F_{S,d}(x,y):=f_S(x).
\]
Let $\gamma_{d-1}$ denote the standard Gaussian probability measure on
$\R^{d-1}$. Then
\[
  \rho_{\Phi_{S,d}}=\rho_S\otimes\gamma_{d-1},
\]
and
\[
  D^2\Phi_{S,d}(x,y)=\diag(\phi''_S(x),I_{d-1}).
\]
Hence
\[
  \nabla F_{S,d}(x,y)=(f'_S(x),0),
  \qquad
  \langle (D^2\Phi_{S,d})^{-1}\nabla F_{S,d},\nabla F_{S,d}\rangle
  =\frac{(f'_S(x))^2}{\phi''_S(x)}.
\]
By Fubini,
\begin{equation}\label{eq:prod-energy}
  \mathcal E_{\Phi_{S,d}}(F_{S,d})=\mathcal E_{\phi_S}(f_S),
\end{equation}
and likewise
\begin{equation}\label{eq:prod-var}
  \Var_{\rho_{\Phi_{S,d}}}(F_{S,d})=\Var_{\rho_S}(f_S),
  \qquad
  \delta_{\BL,\Phi_{S,d}}(F_{S,d})=\delta_{\BL,\phi_S}(f_S).
\end{equation}

The optimizer space is
\[
  \cO_{\BL,\Phi_{S,d}}
  =
  \spanop\{1,\phi'_S(x),y_1,\dots,y_{d-1}\}.
\]
Since $\int f_S\,d\rho_S=0$ and each Gaussian coordinate has zero mean,
$F_{S,d}$ is orthogonal to every generator of this space:
\[
  \int F_{S,d}\,d(\rho_S\otimes\gamma_{d-1})=0,
\]
\[
  \int F_{S,d}(x,y)\phi'_S(x)\,d(\rho_S\otimes\gamma_{d-1})
  =
  \left(\int f_S(x)\phi'_S(x)\,d\rho_S(x)\right)
  \left(\int d\gamma_{d-1}\right)
  =0,
\]
\[
  \int F_{S,d}(x,y)y_j\,d(\rho_S\otimes\gamma_{d-1})
  =
  \left(\int f_S(x)\,d\rho_S(x)\right)
  \left(\int y_j\,d\gamma_{d-1}(y)\right)
  =0.
\]
Therefore
\begin{equation}\label{eq:prod-dist}
  \dist_{L^2(\rho_{\Phi_{S,d}})}(F_{S,d},\cO_{\BL,\Phi_{S,d}})^2
  =
  \|F_{S,d}\|_{L^2(\rho_{\Phi_{S,d}})}^2
  =
  \|f_S\|_{L^2(\rho_S)}^2.
\end{equation}

\begin{theorem}\label{thm:product}
For every integer $d\ge 2$,
\[
  \frac{\dist_{L^2(\rho_{\Phi_{S,d}})}(F_{S,d},\cO_{\BL,\Phi_{S,d}})^2}
       {\delta_{\BL,\Phi_{S,d}}(F_{S,d})}
  =
  S\bigl(1+O(S^{-1})\bigr),
  \qquad S\to\infty.
\]
In particular, for every finite $C$ there exists $S$ such that
$\delta_{\BL,\Phi_{S,d}}(F_{S,d})>0$ and
\[
  \dist_{L^2(\rho_{\Phi_{S,d}})}(F_{S,d},\cO_{\BL,\Phi_{S,d}})^2
  >
  C^2\,\delta_{\BL,\Phi_{S,d}}(F_{S,d}).
\]
\end{theorem}

\begin{proof}
By \eqref{eq:prod-energy}, \eqref{eq:prod-var}, and \eqref{eq:prod-dist}, the
distance squared and deficit for $(\Phi_{S,d},F_{S,d})$ are exactly the
one-dimensional ones for $(\phi_S,f_S)$. The claim therefore follows from
Theorem~\ref{thm:1d}.
\end{proof}

\begin{proof}[Proof of Theorem~\ref{thm:main}]
For $d=1$ this is exactly Theorem~\ref{thm:1d}. For $d\ge 2$ it is
Theorem~\ref{thm:product}. The equivalent formulation with a universal bound
for all $S$ is just the negation of the pointwise contradiction statement.
\end{proof}

\begin{remark}
The argument above proves slightly more than the bare non-existence of a finite
constant: it shows that along the concrete family the ratio
\[
  \frac{\dist_{L^2}^2}{\delta_{\BL}}
\]
actually grows linearly in $S$. The Lean theorem packages only the exact
non-existence statement from Theorem~\ref{thm:main}, but it is obtained from
this stronger asymptotic behavior.
\end{remark}

\section{Intermediate \texorpdfstring{$L^p$}{Lp} distances}

The preceding sections use the Hilbert space distance because the
Brascamp--Lieb variance and the optimizer space are naturally expressed in
$L^2$. The same construction, however, also rules out the direct $L^p$
analogue of the square-root stability estimate for every fixed $1<p<2$.
The only point to check is the size of the distance to the optimizer space in
$L^p$.

For a probability measure $\nu$ and a linear space $V\subset L^p(\nu)$, we use
the notation
\[
  \dist_{L^p(\nu)}(h,V)
  :=
  \inf_{v\in V}\|h-v\|_{L^p(\nu)}.
\]
In the present one-dimensional example,
\[
  \cO_{\BL,\phi_S}=\spanop\{1,\phi'_S\}.
\]
Thus an element of the optimizer space has the form $a\phi'_S+b$.

\begin{lemma}\label{lem:lp-distance-1d}
Let $1\le p<\infty$ be fixed. Along the one-dimensional family constructed
above,
\[
  \dist_{L^p(\rho_S)}(f_S,\cO_{\BL,\phi_S})^p
  \asymp q_S,
  \qquad S\to\infty.
\]
In particular, since $q_S\asymp S^{-1}$,
\[
  \dist_{L^p(\rho_S)}(f_S,\cO_{\BL,\phi_S})
  \asymp S^{-1/p}.
\]
\end{lemma}

\begin{proof}
Recall that
\[
  f_S=g_S-c_S,
  \qquad
  c_S=\int_{\R}g_S\,d\rho_S.
\]
From Lemma~\ref{lem:normalization} and the proof of
Lemma~\ref{lem:variance}, we have
\[
  c_S=q_S+O(t_S)=q_S+O(S^{-3}),
  \qquad
  q_S\asymp S^{-1},
  \qquad
  t_S=O(S^{-3}).
\]
Moreover $g_S=0$ on $C_S$, $g_S=1$ on $E_S$, and $0\le g_S\le 1$ on $T_S$.
Therefore
\[
  f_S=-c_S \quad\text{on } C_S,
  \qquad
  f_S=1-c_S \quad\text{on } E_S,
  \qquad
  |f_S|\le 1+c_S \quad\text{on } T_S.
\]
Since $\rho_S(C_S)=1-q_S-t_S=1+O(S^{-1})$, the function is very small on most
of the mass and of order one on a set of mass $q_S$.

We first remove the odd optimizer direction. For every $a,b\in\R$, the
function $f_S-b$ is even, while $\phi'_S$ is odd and $\rho_S$ is even. Pairing
the points $x$ and $-x$ and using the convexity of $u\mapsto |u|^p$ gives
\[
  \frac{|f_S(x)-b-a\phi'_S(x)|^p
        +|f_S(x)-b+a\phi'_S(x)|^p}{2}
  \ge |f_S(x)-b|^p.
\]
After integration,
\begin{equation}\label{eq:lp-drop-odd}
  \|f_S-a\phi'_S-b\|_{L^p(\rho_S)}^p
  \ge
  \|f_S-b\|_{L^p(\rho_S)}^p.
\end{equation}
Thus the lower bound for the full optimizer space reduces to a lower bound
against constants.

Fix $b\in\R$. If $|b+c_S|\le 1/2$, then on $E_S$,
\[
  |f_S-b|=|1-c_S-b|=|1-(b+c_S)|\ge \frac12,
\]
and hence
\[
  \|f_S-b\|_{L^p(\rho_S)}^p\ge 2^{-p}q_S.
\]
If instead $|b+c_S|>1/2$, then on $C_S$,
\[
  |f_S-b|=|-c_S-b|>\frac12,
\]
so
\[
  \|f_S-b\|_{L^p(\rho_S)}^p
  \ge 2^{-p}\rho_S(C_S).
\]
For all sufficiently large $S$, $\rho_S(C_S)\ge 1/2$ and $q_S\le 1$, so this
second lower bound is also bounded below by a constant multiple of $q_S$.
Together with \eqref{eq:lp-drop-odd}, this proves
\[
  \dist_{L^p(\rho_S)}(f_S,\cO_{\BL,\phi_S})^p\gtrsim q_S.
\]

For the reverse inequality, take the optimizer element to be $0$. Using the
pointwise description above,
\[
  \|f_S\|_{L^p(\rho_S)}^p
  \le
  c_S^p\rho_S(C_S)+(1-c_S)^pq_S+(1+c_S)^pt_S.
\]
Here $c_S\asymp q_S$, $0<q_S<1$, and $t_S=O(S^{-3})=O(q_S)$, so the right-hand
side is $O(q_S)$. Hence
\[
  \dist_{L^p(\rho_S)}(f_S,\cO_{\BL,\phi_S})^p
  \le
  \|f_S\|_{L^p(\rho_S)}^p
  \lesssim q_S.
\]
The two estimates prove the lemma.
\end{proof}

\begin{theorem}\label{thm:lp-failure-1d}
Let $1<p<2$ be fixed. Along the one-dimensional family constructed above,
\[
  \frac{\dist_{L^p(\rho_S)}(f_S,\cO_{\BL,\phi_S})}
       {\delta_{\BL,\phi_S}(f_S)^{1/2}}
  \asymp S^{1-1/p},
\]
and in particular the ratio tends to $+\infty$ as $S\to\infty$. Equivalently,
\[
  \frac{\dist_{L^p(\rho_S)}(f_S,\cO_{\BL,\phi_S})^2}
       {\delta_{\BL,\phi_S}(f_S)}
  \asymp S^{2-2/p}\to+\infty.
\]
Consequently there is no finite constant $C$ such that
\[
  \dist_{L^p(\rho_\phi)}(f,\cO_{\BL,\phi})
  \le C\,\delta_{\BL,\phi}(f)^{1/2}
\]
holds simultaneously for all one-dimensional admissible data containing the
family $(\phi_S,f_S)$.
\end{theorem}

\begin{proof}
By Lemma~\ref{lem:lp-distance-1d},
\[
  \dist_{L^p(\rho_S)}(f_S,\cO_{\BL,\phi_S})
  \asymp S^{-1/p}.
\]
By Proposition~\ref{prop:asymptotics},
\[
  \delta_{\BL,\phi_S}(f_S)^{1/2}
  \asymp S^{-1}.
\]
Therefore
\[
  \frac{\dist_{L^p(\rho_S)}(f_S,\cO_{\BL,\phi_S})}
       {\delta_{\BL,\phi_S}(f_S)^{1/2}}
  \asymp
  \frac{S^{-1/p}}{S^{-1}}
  =
  S^{1-1/p}.
\]
Since $p>1$, the exponent $1-1/p$ is positive, and the ratio diverges. Squaring
the numerator gives the equivalent formulation with exponent $2-2/p$.
\end{proof}

\begin{theorem}\label{thm:lp-product}
Let $d\ge 2$ and $1<p<2$ be fixed. Along the product family
$(\Phi_{S,d},F_{S,d})$,
\[
  \frac{\dist_{L^p(\rho_{\Phi_{S,d}})}
       (F_{S,d},\cO_{\BL,\Phi_{S,d}})}
       {\delta_{\BL,\Phi_{S,d}}(F_{S,d})^{1/2}}
  \asymp S^{1-1/p},
  \qquad S\to\infty.
\]
In particular, for every finite $C$ there exists $S$ such that
$\delta_{\BL,\Phi_{S,d}}(F_{S,d})>0$ and
\[
  \dist_{L^p(\rho_{\Phi_{S,d}})}
       (F_{S,d},\cO_{\BL,\Phi_{S,d}})
  >
  C\,\delta_{\BL,\Phi_{S,d}}(F_{S,d})^{1/2}.
\]
\end{theorem}

\begin{proof}
As in the $L^2$ product argument,
\[
  \cO_{\BL,\Phi_{S,d}}
  =
  \spanop\{1,\phi'_S(x),y_1,\dots,y_{d-1}\}.
\]
Let
\[
  H(x,y):=F_{S,d}(x,y)-b=f_S(x)-b.
\]
This function is even in $x$ and invariant under every sign change
$y_j\mapsto -y_j$. The probability measure
$\rho_{\Phi_{S,d}}=\rho_S\otimes\gamma_{d-1}$ has the same symmetries, while
$\phi'_S(x)$ is odd in $x$ and each $y_j$ is odd in its own Gaussian variable.
Therefore, by averaging over the finite group generated by these sign changes
and using convexity of $u\mapsto |u|^p$, for every
$a\in\R$, $\alpha\in\R^{d-1}$, and $b\in\R$,
\[
  \|F_{S,d}-a\phi'_S-\alpha\cdot y-b\|_{L^p(\rho_{\Phi_{S,d}})}^p
  \ge
  \|F_{S,d}-b\|_{L^p(\rho_{\Phi_{S,d}})}^p.
\]
Since $F_{S,d}$ is independent of $y$, the right-hand side equals
\[
  \|f_S-b\|_{L^p(\rho_S)}^p.
\]
The lower-bound argument from Lemma~\ref{lem:lp-distance-1d}, applied to
constants $b$, gives
\[
  \dist_{L^p(\rho_{\Phi_{S,d}})}
       (F_{S,d},\cO_{\BL,\Phi_{S,d}})^p
  \gtrsim q_S.
\]
Conversely, taking the zero optimizer gives
\[
  \dist_{L^p(\rho_{\Phi_{S,d}})}
       (F_{S,d},\cO_{\BL,\Phi_{S,d}})^p
  \le
  \|F_{S,d}\|_{L^p(\rho_{\Phi_{S,d}})}^p
  =
  \|f_S\|_{L^p(\rho_S)}^p
  \lesssim q_S.
\]
Thus the product-family $L^p$ distance has the same size as the
one-dimensional one:
\[
  \dist_{L^p(\rho_{\Phi_{S,d}})}
       (F_{S,d},\cO_{\BL,\Phi_{S,d}})
  \asymp S^{-1/p}.
\]
The deficit is unchanged by the product lift, by \eqref{eq:prod-var}, so
\[
  \delta_{\BL,\Phi_{S,d}}(F_{S,d})^{1/2}
  =
  \delta_{\BL,\phi_S}(f_S)^{1/2}
  \asymp S^{-1}.
\]
Dividing gives the claimed asymptotic ratio.
\end{proof}

\begin{remark}
At $p=1$ the same computation is exactly borderline. Lemma~\ref{lem:lp-distance-1d}
gives
\[
  \dist_{L^1(\rho_S)}(f_S,\cO_{\BL,\phi_S})
  \asymp S^{-1},
\]
while Proposition~\ref{prop:asymptotics} gives
\[
  \delta_{\BL,\phi_S}(f_S)^{1/2}\asymp S^{-1}.
\]
Thus this counterexample does not contradict an $L^1$ square-root stability
estimate. In fact, the companion paper \cite{BLpaper} proves precisely such a
sharp uniform $L^1$ stability theorem for the Brascamp--Lieb deficit. The
present construction shows that the topology in that theorem cannot be
strengthened to any fixed $L^p$ topology with $p>1$ by a dimension-free
constant of the same square-root form.
\end{remark}

\begin{thebibliography}{10}

\bibitem{BLpaper}
J.~M. Machado and J.~P. G. Ramos,
\textit{Quantitative stability for Brascamp--Lieb and moment measures},
arXiv:2511.22636v2.

\bibitem{BL76}
H.~J. Brascamp and E.~H. Lieb,
\textit{On extensions of the Brunn--Minkowski and Pr\'ekopa--Leindler theorems,
including inequalities for log-concave functions, and with an application to
the diffusion equation},
J.~Funct.~Anal. 22 (1976), 366--389.

\end{thebibliography}

\end{document}
